\chapter{CONCLUSION}

\hspace{0.5cm} The proposed AI-Powered Code Document Generator would thus offer an intelligent and automated solution for the generation of accurate, structured, and context-aware documentation from source code. This would thus help in the identification of significant software components like functions, classes, imports, call relationships, entry points, and recursive patterns, among others. In this regard, the proposed solution would leverage Tree-sitter for parsing code across fourteen programming languages and thus generate accurate documentation like the README file, architecture, description of files, recent changes of the git repository, among others.

\par The system relies on a three-stage feature extraction process to guarantee the LLM access to rich and non-redundant information for every function it documents. Features extracted from the AST offer structural information such as the function’s name, class membership, parameters, return type, and calls. Features extracted from the text offer context-based information on the file’s complexity and documentation. Features extracted from the graph offer architectural information on the function and the program based on the NetworkX library’s directed call graph. By combining all the features in a single fact block, the system offers documentation that is functionally correct, contextually scaled and architecturally aware.

\par This modular and language-independent design ensures that each module is independent in its functioning through well-defined interfaces. It is possible to upgrade or extend any module in the pipeline without affecting the rest of the pipeline. Cross-file awareness ensures that functions used across module boundaries are correctly identified as active rather than unused, thereby improving the reliability of the architectural analysis. The parallel document generation using the \texttt{ThreadPoolExecutor} ensures the pipeline scales well with projects of different sizes.

\par This tool is highly beneficial for software development teams working with repositories that lack adequate documentation. It boosts productivity by automating what is usually a manual and time-consuming process. Additionally, it supports long-term code maintainability by allowing developers to understand the code’s structure through graph-based analysis.

\par Overall, the project highlights the potential of combining artificial intelligence with static code analysis to create a practical tool that simplifies software documentation. Future improvements could include extending support to more programming languages, integrating with version control systems to enable real-time documentation updates and employing more advanced language models for richer, more detailed explanations, ultimately developing the system into a comprehensive, production-ready documentation platform.
\newline
\par Despite the current system addressing many key challenges in automated documentation, there is room for improvement by enhancing system capabilities, supporting additional programming languages and adding new features to make the system more intelligent and adaptive. These opportunities pave the way for future developments that could increase the system’s effectiveness and applicability in real-world software development scenarios which discussed in next chapter.